\begin{pastexam}[label={px:ncu-107-dsa-14}]{中央資工所 107 學年度 第 14 題}
（單選；每題 5 分，答錯倒扣一分）

A hash function maps a key into a bucket in the hash table. Which of the following statements is
false?
\begin{enumerate}[label=(\alph*)]
\item A Biased use of the hash table is not a desired property of hash functions.
\item A division hash function with divisor $D=2^r$, where $r$ is an integer, may result in
serious collision.
\item Doubling hash table (i.e., rebuild the hash table by doubling the number of table buckets)
may result in key search suspended for unacceptable long periods while the rebuild is in progress.
\item Dynamic hash reduces the table rebuild time by ensuring a small amount of rebuild changes
\item When chaining is used to resolve overflows, the search for a key involves comparison with
keys that have different hash values.
\end{enumerate}
\end{pastexam}

\begin{pastsoln}
答案是 (e)。 題目問哪一個是\textbf{錯}的。

\begin{itemize}
\item (a) 對。 雜湊函數要把鍵盡量平均分到每個桶（資結書第~12.1~節的簡單均勻雜湊）。 偏向某些桶
（biased use）會讓那些桶的串列變長或探測變多，搜尋變慢，不是想要的性質。
\item (b) 對。 除法 $h(k) = k \bmod 2^r$ 只取 $k$ 最低的 $r$ 個位元，高位元完全沒用到（資結書第~12.1~節）。
例如鍵都是 $2^r$ 的倍數時，全部落在第 $0$ 個桶；低位元規律相同的鍵（例如都是偶數、結尾相同的編號）
也會擠在少數幾個桶裡。
\item (c) 對。 表加倍（再雜湊，資結書第~12.5~節）時要配置新表，把舊表的每一個鍵重新算雜湊值、搬進新表，
一次要 $\Theta(n + m)$。 平攤下來每次插入是 $O(1)$，但搬的那一次很久；表很大或在磁碟上時，搬動期間搜尋
得等，暫停時間可能長到不能接受。
\item (d) 對。 動態雜湊（例如資結書第~12.5~節的可延伸雜湊）一次只拆一個桶、只搬那個桶裡的鍵，目錄加倍時也
只複製指標，每次重建只改動一小部分，這正是它要解決 (c) 的問題。
\item (e) \textbf{錯}。 鏈結法（資結書第~12.2~節）把雜湊值相同的鍵放在同一條串列，搜尋鍵 $k$ 時只走
$T[h(k)]$ 那一條，比較的都是\textbf{雜湊值和 $k$ 相同}的鍵，不會和雜湊值不同的鍵比較。 會和不同雜湊值的
鍵比較的是開放定址，例如線性探測時 $k$ 會一路比過被別的鍵佔走的格子（資結書第~12.3~節）。
\end{itemize}
\end{pastsoln}
